\documentclass[12pt]{article}
\usepackage[utf8]{inputenc}
\usepackage[T1]{fontenc}
\usepackage[french]{babel}
\usepackage{booktabs}
\usepackage{pgfplots}
\usepackage{graphicx}
\pgfplotsset{compat=1.18}

\title{Benchmark des piles génériques : tableau vs liste chaînée}
\author{\textbf{[NOM Prénom]}}
\date{\today}

\begin{document}
\maketitle

\section{Introduction}
Ce rapport présente les performances de deux implémentations d'une pile générique en C :
une basée sur un tableau dynamique et une autre sur une liste chaînée simple.
Nous mesurons le temps d'exécution pour plusieurs opérations (push, pop, mixte) et
différentes tailles de données (int, double, structure de 64~octets).
Nous étudions également l'impact de la mise en œuvre sous forme de DLL (liaison
implicite et explicite) par rapport au lien statique.

\section{Implémentations}
\begin{description}
  \item[Tableau dynamique] \hfill\\
    La pile est un bloc mémoire contigu qui double sa capacité lorsqu'il est plein.
    L'amortissement donne un coût moyen \(\mathcal{O}(1)\) par \texttt{push}, mais
    nécessite occasionnellement un \texttt{realloc} coûteux.
  \item[Liste chaînée] \hfill\\
    Chaque élément est une allocation séparée (\texttt{malloc}/\texttt{free}).
    Pas de redimensionnement, mais chaque opération implique un appel au
    allocateur et une moins bonne localité de cache.
\end{description}

\section{Méthodologie de benchmarking}
\begin{itemize}
  \item Chronométrage haute précision avec \texttt{QueryPerformanceCounter}.
  \item Chaque configuration répétée 10 fois (pour les tailles testées) ; on rapporte moyenne, minimum et
        écart-type.
  \item Tailles de poussée : \(N = 10^3, 10^4\).
  \item Types d'éléments : \texttt{int} (4~B), \texttt{double} (8~B), struct de
        64~B.
  \item Quatre scénarios :
        \begin{enumerate}
          \item uniquement des \texttt{push},
          \item \texttt{push} suivi de \texttt{pop} (pile pleine puis vide),
          \item alternance \texttt{push}/\texttt{pop},
          \item \texttt{N pushes} puis \texttt{N pops}.
        \end{enumerate}
  \item Le code est compilé avec \texttt{-O2} qui active les optimisations ; une somme volatile des valeurs dépilées est utilisée pour
        empêcher l'élimination du code mort.
\end{itemize}

\section{Résultats}
Les chiffres suivants proviennent des fichiers \texttt{bench\_tableau.csv} et
\texttt{bench\_liste.csv} générés par les exécutables \texttt{bench\_tableau.exe}
et \texttt{bench\_liste.exe}. Les temps sont exprimés en millisecondes (ms).

\subsection{Tableau dynamique}
\begin{table}[ht]
\centering
\caption{Temps moyen (en ms) pour \texttt{N pushes}, type \texttt{int}.}
\begin{tabular}{lrr}
\toprule
\(N\) & \(10^3\) & \(10^4\) \\
\midrule
Temps & 0.069750 & 0.787400 \\
\bottomrule
\end{tabular}
\label{tab:tableau-int-push}
\end{table}

\subsection{Liste chaînée}
\begin{table}[ht]
\centering
\caption{Temps moyen (en ms) pour \texttt{N pushes}, type \texttt{int}.}
\begin{tabular}{lrr}
\toprule
\(N\) & \(10^3\) & \(10^4\) \\
\midrule
Temps & 0.623680 & 7.264290 \\
\bottomrule
\end{tabular}
\label{tab:liste-int-push}
\end{table}

\subsection{Ratio L/T (temps liste / temps tableau)}
\begin{table}[ht]
\centering
\caption{Ratio du temps moyen de la liste chaînée par rapport au tableau dynamique pour \texttt{N pushes}, type \texttt{int}.}
\begin{tabular}{lrr}
\toprule
\(N\) & \(10^3\) & \(10^4\) \\
\midrule
Ratio L/T & 8.94 & 9.23 \\
\bottomrule
\end{tabular}
\label{tab:ratio-lt}
\end{table}

\subsection{Graphique de comparaison (push uniquement)}
\begin{figure}[ht]
\centering
\begin{tikzpicture}
\begin{axis}[
    xlabel={Nombre d'opérations $N$},
    ylabel={Temps moyen [ms]},
    xmode=log,
    ymode=log,
    legend pos=north west,
    grid=both,
    width=0.8\textwidth,
    height=0.5\textwidth
]
\addplot+[mark=square,blue] table[x index=0, y index=1]{
N temps
1e3 0.069750
1e4 0.787400
};
\addlegendentry{Tableau}
\addplot+[mark=triangle,red] table[x index=0, y index=1]{
N temps
1e3 0.623680
1e4 7.264290
};
\addlegendentry{Liste chaînée}
\end{axis}
\end{tikzpicture}
\caption{Temps moyen pour des \texttt{N pushes} (type \texttt{int}) en fonction de $N$ (en ms).}
\label{fig:push-int}
\end{figure}

\section{Comparaison avant DLL vs après DLL}
Nous avons exécuté le même benchmark sous trois configurations :
lien statique, DLL avec liaison implicite, et DLL avec liaison explicite.
Cependant, seules les mesures pour la liste chaînée en liaison implicite sont disponibles (voir tableau~\ref{tab:dll-overhead}).
Les autres combinaisons (tableau en liaison implicite/explicite, liste en liaison explicite) manquaient lors de la rédaction de ce rapport.
Le tableau~\ref{tab:dll-overhead} montre le surcoût (en \%) de la liaison implicite par rapport au lien statique pour la liste chaînée, type \texttt{int}, opération \texttt{N pushes}.
Les valeurs négatives indiquent une amélioration (probablement due à la variabilité des mesures).

\begin{table}[ht]
\centering
\caption{Surcoût de la liaison implicite DLL par rapport au lien statique (liste chaînée, type \texttt{int}, \texttt{N pushes}).}
\begin{tabular}{lrr}
\toprule
\(N\) & \(10^3\) & \(10^4\) \\
\midrule
Surcoût (\%) & -17.01 & 25.96 \\
\bottomrule
\end{tabular}
\label{tab:dll-overhead}
\end{table}

Le surcoût observé pour la liaison explicite provient de l'appel indirect via pointeur de fonction ajouté à chaque appel de fonction de la pile.
Le chargement de la DLL et la résolution des symboles avec \texttt{GetProcAddress} ne sont effectués qu'une seule fois au démarrage,
donc leur impact est négligeable dans les benchmarks de longue durée.
La liaison implicite utilise également un appel indirect, mais la résolution de l'adresse de la fonction est effectuée au chargement
via la table d'import (IAT), ce qui évite l'appel à \texttt{GetProcAddress} lors de chaque utilisation.

\section{Analyse et justification}
\begin{itemize}
  \item \textbf{Localité de cache} : le tableau bénéficie d'un accès séquentiel
        et d'un taux de hit élevé dans les caches L1/L2, tandis que chaque nœud
        de la liste provoque un chargement de ligne de cache distinct.
  \item \textbf{Coût d'allocation} : la liste effectue un \texttt{malloc} et un
        \texttt{free} par élément, tandis que le tableau n'alloue que lors
        des augmentations de capacité (coût amorti \(\mathcal{O}(1)\)).
  \item \textbf{Impact de la DLL} : le surcoût provient surtout du passage par
        la table d'import (\texttt{__imp\_}) pour la liaison implicite, ou de
        l'appel indirect via pointeur de fonction pour la liaison explicite.
        Dans les deux cas, l'impact est minimal pour les charges de travail longues.
\end{itemize}

\section{Conclusion et recommandations}
\begin{itemize}
  \item Pour toutes les tailles de pile testées (\(N = 10^3\) et \(10^4\)), le
        \textbf{tableau dynamique} est nettement plus rapide que la liste chaînée
        grâce à sa meilleure localité de cache et à son moindre nombre d'appels
        au allocateur (voir tableau~\ref{tab:ratio-lt}).
  \item L'utilisation sous forme de DLL n'entraîne pas de pénalité significative
        pour les charges de travail longues ; choisir la liaison implicite
        pour sa simplicité, sauf si le chargement différé ou le changement
        d'implémentation à l'exécution est requis.
\end{itemize}

\section{Environnement de test}
\begin{itemize}
  \item Processeur : Intel Core i5-8250U
  \item Mémoire : 8 Go de RAM DDR4
  \item Système d'exploitation : Windows 11
  \item Compilateur : \textbf{[À COMPLÉTER : version de gcc]} (MinGW-w64)
  \item Options de compilation : \texttt{-O2 -Wall -Wextra} avec
        \texttt{-DPILE\_STATIC} pour les versions statiques, et
        \texttt{-DPILE\_EXPORTS} lors de la compilation des DLL.
\end{itemize}

\section{Guide d'exécution pour l'enseignant}
Ce guide explique comment compiler, exécuter et tester l'ensemble du benchmark.

\subsection{Prérequis}
\begin{itemize}
  \item MinGW-w64 installé avec \texttt{gcc} dans le \texttt{PATH}.
  \item PowerShell (invite de commandes Windows fonctionne également).
  \item Windows 64 bits (les exécutables générés sont 64 bits).
\end{itemize}
Pour vérifier l'installation de gcc, ouvrez un terminal et exécutez :
\begin{verbatim>
gcc --version
\end{verbatim}
Vous devriez voir quelque chose comme :
\begin{verbatim>
gcc.exe (x86_64-posix-seh-rev0, Built by MinGW-WinBuild project) 11.2.0
\end{verbatim}
(le numéro de version peut varier).

\subsection{Compilation}
Ouvrez un terminal dans le dossier du projet et exécutez :
\begin{verbatim>
.\build.bat
\end{verbatim}
Le script générera les fichiers suivants :
\begin{itemize}
  \item \texttt{bench\_tableau.exe} et \texttt{bench\_liste.exe} : versions statiques du benchmark.
  \item \texttt{pile\_tableau.dll} et \texttt{pile\_liste.dll} : les bibliothèques dynamiques.
  \item \texttt{libpile\_tableau.dll.a} et \texttt{libpile\_liste.dll.a} : bibliothèques d'import pour la liaison implicite.
  \item \texttt{bench\_dll\_tableau\_implicit.exe} et \texttt{bench\_dll\_liste\_implicit.exe} : versions du benchmark liées implicitement aux DLL.
  \item \texttt{bench\_dll.exe} : version du benchmark utilisant la liaison explicite (chargement exécutif).
  \item \texttt{client\_implicit\_tableau.exe} et \texttt{client\_implicit\_liste.exe} : clients utilisant la liaison implicite.
  \item \texttt{client\_explicit.exe} : client utilisant la liaison explicite.
\end intermédiaire objets (\texttt{*.o}) sont supprimés à la fin du script.

\subsection{Exécution des benchmarks de base}
Les exécutables statiques produisent directement des fichiers CSV :
\begin{verbatim>
.\bench_tableau.exe > bench_tableau.csv
.\bench_liste.exe > bench_liste.csv
\end{verbatim}
Les fichiers \texttt{bench\_tableau.csv} et \texttt{bench\_liste.csv} contiennent les résultats pour les opérations de push et de pop
avec les tailles \(N = 10^3\) et \(10^4\) et le type \texttt{int}.

\subsection{Exécution des tests DLL}
\subsubsection{Liaison implicite (versions pré-liées)}
Les exécutables suivants sont déjà liés aux DLL et peuvent être exécutés directement :
\begin{verbatim>
.\bench_dll_tableau_implicit.exe > bench_dll_tableau_implicit.csv
.\bench_dll_liste_implicit.exe > bench_dll_liste_implicit.csv
\end{verbatim}

\subsubsection{Liaison explicite (chargement exécutif)}
L'exécutable \texttt{bench\_dll.exe} nécessite deux arguments : l'implémentation (\texttt{tableau} ou \texttt{liste}) et le mode (\texttt{static}, \texttt{implicit} ou \texttt{explicit}).
Cependant, le mode \texttt{static} n'est pas pertinent ici puisqu'il s'agirait de nouveau de la liaison statique.
Pour tester la DLL avec liaison explicite, utilisez :
\begin{verbatim>
.\bench_dll.exe tableau explicit > bench_dll_tableau_explicit.csv
.\bench_dll.exe liste explicit > bench_dll_liste_explicit.csv
\end{verbatim}
Note : le mode \texttt{implicit} avec \texttt{bench\_dll.exe} effectue un chargement explicite de la DLL mais utilise ensuite la table d'import obtenu via \texttt{GetProcAddress} pour chaque appel ?
En fait, dans notre code, le mode \texttt{implicit} de \texttt{bench\_dll.exe} utilise la liaison implicite via l'import lib, donc il est identique aux exécutables pré-liés.
Nous recommandons d'utiliser les exécutables pré-liés pour la liaison implicite afin de éviter toute confusion.

\subsection{Les clients de test des DLL}
Les fichiers \texttt{client\_implicit.c} et \texttt{client\_explicit.c} illustrent l'utilisation des DLL.
Ils ont déjà été compilés par \texttt{build.bat} :
\begin{verbatim>
.\client_implicit_tableau.exe
.\client_implicit_liste.exe
.\client_explicit.exe
\end{verbatim}
Ces programmes effectuent quelques opérations de pile et affichent un message de réussite.

\subsection{Vérification que les DLL sont bien utilisées}
Pour vérifier qu'un exécutable utilise bien une DLL, vous pouvez lister les dépendances avec \texttt{objdump} :
\begin{verbatim>
objdump -p mon_executable.exe | Select-String "DLL Name"
\end{verbatim}
Sous PowerShell, cela affiche les noms des DLL importées.
Pour lister les exportations d'une DLL, utilisez :
\begin{verbatim>
objdump -p pile_liste.dll
\end{verbatim}
Recherchez la section contenant les symboles exportés.

\subsection{Remarque importante}
Les fichiers \texttt{*.dll} doivent se trouver dans le même répertoire que l'exécutable qui les utilise, sinon le chargement échouera.
Assurez-vous de lancer les exécutables depuis le dossier du projet où les DLL ont été générées.

\subsection{Conseil d'affichage}
Dans PowerShell, les caractères accentués peuvent ne pas s'afficher correctement.
Pour corriger cela, changez la page de code en UTF-8 avant d'exécuter les commandes :
\begin{verbatim>
chcp 65001
\end{verbatim}

\subsection{Erreurs fréquentes et leurs solutions}
\begin{description}
  \item[DLL introuvable] :
        Assurez-vous que la DLL est dans le même dossier que l'exécutable ou dans un répertoire référencé par le \texttt{PATH}.
  \item[Erreur 193] :
        Cette erreur indique un mélange entre binaire 32 bits et 64 bits.
        Vérifiez que vous utilisez bien une version 64 bits de MinGW-w64 et que vous n'essayez pas d'exécuter un binaire 32 bits sur un système 64 bits (ou vice versa).
  \item[Fonction introuvable avec la liaison explicite] :
        Vérifiez que le nom de la fonction passé à \texttt{GetProcAddress} est exactement celui exporté par la DLL (veillez à la casse et à l'absence de décoration).
        Vous pouvez lister les exportations réelles avec \texttt{objdump -p <dll>}.
\end{description}

\end{document}